% HD-S28. GPT-6 Astra Pro, /root/sixfolds, inherited effort not exposed.
% 21 November 2026. Exact scope of the full natural E8 scalar-Weyl route.
\subsection{HD-S28: rank-nine natural \(E_8\) sixfolds have determinant \([-1]\) or \([-3]\)}
\label{sec:e8-determinant-scope}

For the full group \(G=W(E_8)\cdot\mu_4\), every natural
primitive Weil sixfold of signature \((3,3)\) arising from a
connected sphere cover with at least twelve true branch values has
ordinary Hermitian determinant class \([-1]\) or \([-3]\).
Every such cover with fewer branch values has period rank at most
eight. In particular this full natural scalar-Weyl construction
cannot yield a rank-nine family in any other determinant component.
The assertion concerns the actual restricted-theta polarization of
the primitive factor of HD-S22, and the specified polarized lattice
transports of HD-S24.

\paragraph{1. A primitive seventh root cannot be a local eigenvalue.}
Let \((g_1,\ldots,g_b)\) be the generating product-one tuple of
nonidentity branch elements, and let \(c_j=\operatorname{rank}(1-g_j)\).
The exact parabolic quotient gives \(\sum_jc_j=22\). There are
at least two entries outside \(W(E_8)\), and each has
\(c_j\ge4\), by the conjugate \(\pm i\) eigenspace argument
in HD-S26. For \(b\ge12\) the total excess satisfies
\[
 \sum_j(c_j-1)=22-b\le10.
\]
The two scalar-odd entries already contribute at least six to this
excess. A real entry with \(c_j\ge6\) would add at least five,
which is impossible. Thus every real branch element has
\(c_j\le5\).

If a real branch element had a primitive seventh root as an
eigenvalue, its rational characteristic polynomial would contain
\(\Phi_7\), of degree six. Those six eigenvalues are all different
from one, contradicting \(c_j\le5\). If a scalar-odd branch
element \(iw\) had such an eigenvalue \(\zeta_7^a\), the
real rational matrix \(w\) would have eigenvalue
\(-i\zeta_7^a\), of exact order 28. Its characteristic polynomial
would then contain \(\Phi_{28}\), of degree 12, impossible for
an \(8\times8\) matrix. Hence no branch element has a primitive
seventh-root eigenvalue.

\paragraph{2. An exact integral invariant projector.}
The previously computed group order factors as
\[
 |G|=1393459200=2^{15}3^5 5^2 7.
\]
For each local matrix \(g\), write its Gaussian characteristic
polynomial as
\[
 \chi_{K,g}(T)=(T-1)^k f_g(T),\qquad
 f_g\in\mathbb Z[i][T],\quad f_g(1)\ne0.
\]
The rational characteristic polynomial is
\(\chi_{K,g}\bar\chi_{K,g}\). Since \(g\) has finite order,
its nontrivial factors are cyclotomic polynomials \(\Phi_r\),
where \(r\mid |G|\). For \(r>1\),
\(\Phi_r(1)\) equals a prime \(p\) when \(r\) is a power of
\(p\), and equals one otherwise. This follows without a
factorization assumption: evaluating
\((T^r-1)/(T-1)=\prod_{d\mid r,d>1}\Phi_d(T)\) at one and
inducting on \(r\) gives the assertion; for a number with at
least two prime divisors its proper prime-power divisors already
contribute the full product \(r\).

The factor 7 occurs to exponent one in \(|G|\), and paragraph 1
excludes \(\Phi_7\). It follows that
\[
 N_{K/\mathbb Q}\bigl(f_g(1)\bigr)
                         \in\{2^a3^b5^c:a,b,c\ge0\}.
\]
Thus \(f_g(1)\) is a unit in \(R=\mathbb Z[i,1/30]\).
On the Gaussian root lattice over \(R\), define the actual
endomorphism
\[
                 \pi_g=\frac{f_g(g)}{f_g(1)}.
\]
The matrix \(g\) is semisimple over characteristic zero, so this
is exactly its projector onto the fixed space; consequently
\(\pi_g^2=\pi_g\) and \(g\pi_g=\pi_g\) as identities of
integral \(R\)-matrices. Write
\(f_g(T)-f_g(1)=(T-1)q_g(T)\). Then
\[
 (1-g)\frac{q_g(g)}{f_g(1)}=1-\pi_g.
\]
This supplies an inverse to \(1-g\) on the complementary summand.
In particular the maps induced by \(\pi_g\) identify
\[
 L/(g-1)L\ \xrightarrow{\ \sim\ }\ \pi_gL=L^g.
\]
Both the invariant and coinvariant boundary modules are free direct
summands. No division by an inertia order containing 7 has been used.

\paragraph{3. Perfect parabolic form and its determinant.}
The \(E_8\) root Gram matrix has determinant one, so its
coefficient lattice is Hermitian self-dual over \(R\). The central
scalar \(iI\) lies in the generated monodromy group, and
\(i-1\) is a unit in \(R\), so the two absolute zeroth
homology groups for the coefficient and dual systems vanish.

Apply the exact bordered-surface argument of HD-S27, paragraph 2,
over this PID. Its only use of an invertible inertia order was to
make the invariant and coinvariant boundary modules free. The
explicit \(\pi_g\) maps in paragraph 2 give precisely those
same boundary modules and identifications here. Absolute and
relative first homology are free and have the perfect
Poincar\'e--Lefschetz cross-pairing. The image of absolute in
relative homology is primitive, since its cokernel is the free
boundary zeroth homology. Extending functionals from that primitive
image, and using the adjoint absolute-to-relative maps, gives a
perfect parabolic cup pairing. Thus the rank-six Hermitian cup
lattice is self-dual over \(R\).

Its Gram determinant is a rational unit of \(R\), with negative
sign from signature \((3,3)\). Therefore it is
\[
                         -2^a3^b5^c\qquad(a,b,c\in\mathbb Z).
\]
The identities \(2=N(1+i)\) and \(5=N(2+i)\) remove the first
and last factors in \(\mathbb Q^\times/N(K^\times)\); only the
parity of \(b\) remains. The class is \([-1]\) or \([-3]\).
Finally the exact coefficient-to-Jacobian map of HD-S23 multiplies
the Hermitian form by \(D=|G|/8\). Its rank-six determinant
multiplier \(D^6=N(D^3)\) is a norm. The same conclusion therefore
holds for the actual primitive restricted-theta factor, and the
specified isometry or similitude lattice maps preserve it.

\paragraph{4. The order-28 candidate has a stronger split conclusion.}
There are actual generating twelve-point tuples with codimensions
\((4,1^{10},8)\) whose last element is \(v=iw_7\), where
\(\chi_{w_7}(T)=\Phi_7(T)(T-1)^2\). Construct them by taking
\(u=ia^5\), all eight simple reflections and two further root
reflections in a selected order, with product \(P\), and setting
\(w_7=P^{-1}a^5\). The four explicit choices in
\texttt{e8\_order7\_candidates.json} satisfy \(w_7^7=1\), the
stated characteristic polynomial, and \(uPv=1\). All simple
reflections occur, so the tuple generates \(W(E_8)\); adjoining
\(u\) recovers the scalar \(i\), hence gives the full group.
The local ages are \(2,\tfrac12^{\,10},4\), giving signature
\((3,3)\).

For this entire local profile, not just those four tuples, the
nontrivial-eigenvalue factors at one are \(16\) for \(u\),
2 for every reflection, and
\[
 \det(1-iw_7)=(1-i)^2\,i^6\Phi_7(-i)=2
\]
for the last branch. The same projector and parabolic-duality proof
therefore works already over \(\mathbb Z[i,1/2]\). Every cover
with this profile is split. The four direct cup computations give
Gram determinants \(-8,-2,-2,-2\), in agreement with the proof.
They do not assert period rank nine, and no supplier-cycle conclusion
is imported from their parameter count.

\paragraph{Check record and exact scope.}
The proof uses the freshly computed finite group and the explicit
integral lattice, followed by the exact local projector, boundary
maps, and primitive-factor map displayed above. It does not infer a
universal determinant result from sampled tuples. It restricts this
natural full-group sphere-cover construction; no assertion excluding
other groups, representations, geometric constructions or cycle
correspondences is made.
